% ─────────────────────────────────────────────────────────────────────────────
% D-Cache Controller FSM + Scratchpad FSM
% Sources: ip_dcache/src/asDCache.sv, ip_scratch/src/asScratch.sv
% ─────────────────────────────────────────────────────────────────────────────

\paragraph{D-Cache FSM overview.}
The \texttt{asDCache} module implements a write-back, write-allocate
4-way set-associative cache with a full 11-state Moore-FSM
(Fig.~\ref{fig:asDCacheFsm01}).
The three columns in the figure reflect the three independent activity
paths:

\begin{itemize}
  \item \textbf{Centre column (normal cache path):}
    \textsc{Idle} $\to$ \textsc{Lookup} $\to$ \textsc{Miss} $\to$
    \textsc{Fill} $\to$ \textsc{Fill\_Resp} $\to$ \textsc{Idle}.
  \item \textbf{Left column (dirty-eviction path):}
    \textsc{Evict\_AW} $\to$ \textsc{Evict\_W} $\to$ \textsc{Evict\_B};
    entered from \textsc{Lookup} (miss on dirty victim) or from the
    flush path.
  \item \textbf{Right column (flush path):}
    \textsc{Flush} $\to$ \textsc{Flush\_Evict\_Prep} $\to$
    \textsc{Evict\_AW}; entered from \textsc{Idle} on
    \texttt{dc\_flush\_i}.
\end{itemize}

\noindent
The \texttt{ERR} state (reached from \textsc{Fill} and \textsc{Evict\_B}
on AXI4 error response) returns unconditionally to \textsc{Idle}.
Self-loops that keep a state while an AXI4 handshake is pending
(\textsc{Miss}, \textsc{Evict\_AW}, \textsc{Evict\_B}) are shown for
completeness.
$\mathit{flush\_mode\_r}$ is a one-bit register that distinguishes
a normal dirty-victim eviction from a flush-triggered write-back.

\begin{figure}[H]
\centering
\resizebox{0.98\textwidth}{!}{%
\begin{tikzpicture}[shorten >=1pt, node distance=1pt, auto,
  every state/.style={fill=blue!10}, initial text=]

  % ── Absolute positions (cm): three-column layout ─────────────────────────
  % Col A (eviction): x = -6
  % Col B (main):     x =  0
  % Col C (flush):    x =  6
  % Rows:             y = 0, -3, -6, -9, -12

  % ── State nodes ──────────────────────────────────────────────────────────

  % Centre column
  \node[state with output, initial above, accepting] (idle) at (0,0)
    { \texttt{IDLE}
      \nodepart{lower}
      \begin{tiny}$\begin{array}{l}
        \mathit{stall}{=}0
      \end{array}$\end{tiny}
    };

  \node[state with output] (lookup) at (0,-3)
    { \texttt{LOOKUP}
      \nodepart{lower}
      \begin{tiny}$\begin{array}{l}
        \mathit{stall}{=}1
      \end{array}$\end{tiny}
    };

  \node[state with output] (miss) at (0,-6)
    { \texttt{MISS}
      \nodepart{lower}
      \begin{tiny}$\begin{array}{l}
        \mathit{stall}{=}1 \\
        \mathit{arvalid}{=}1
      \end{array}$\end{tiny}
    };

  \node[state with output] (fill) at (0,-9)
    { \texttt{FILL}
      \nodepart{lower}
      \begin{tiny}$\begin{array}{l}
        \mathit{stall}{=}1 \\
        \mathit{rready}{=}1
      \end{array}$\end{tiny}
    };

  \node[state with output] (fill_resp) at (0,-12)
    { \texttt{FILL\_RESP}
      \nodepart{lower}
      \begin{tiny}$\begin{array}{l}
        \mathit{stall}{=}0 \\
        \mathit{rvalid}{=}{\neg}\mathit{wr}
      \end{array}$\end{tiny}
    };

  % Left column (eviction chain)
  \node[state with output] (evict_aw) at (-6,-3)
    { \texttt{EVICT\_AW}
      \nodepart{lower}
      \begin{tiny}$\begin{array}{l}
        \mathit{stall}{=}1 \\
        \mathit{awvalid}{=}1
      \end{array}$\end{tiny}
    };

  \node[state with output] (evict_w) at (-6,-6)
    { \texttt{EVICT\_W}
      \nodepart{lower}
      \begin{tiny}$\begin{array}{l}
        \mathit{stall}{=}1 \\
        \mathit{wvalid}{=}1
      \end{array}$\end{tiny}
    };

  \node[state with output] (evict_b) at (-6,-9)
    { \texttt{EVICT\_B}
      \nodepart{lower}
      \begin{tiny}$\begin{array}{l}
        \mathit{stall}{=}1 \\
        \mathit{bready}{=}1
      \end{array}$\end{tiny}
    };

  % Right column (flush path)
  \node[state with output] (flush) at (6,0)
    { \texttt{FLUSH}
      \nodepart{lower}
      \begin{tiny}$\begin{array}{l}
        \mathit{stall}{=}1
      \end{array}$\end{tiny}
    };

  \node[state with output] (flush_ep) at (6,-3)
    { \texttt{FLUSH\_EVICT\_PREP}
      \nodepart{lower}
      \begin{tiny}$\begin{array}{l}
        \mathit{stall}{=}1
      \end{array}$\end{tiny}
    };

  % ERR (right of FILL)
  \node[state with output] (err) at (5,-9)
    { \texttt{ERR}
      \nodepart{lower}
      \begin{tiny}$\begin{array}{l}
        \mathit{err}{=}1 \\
        \mathit{stall}{=}0
      \end{array}$\end{tiny}
    };

  % ── Transitions ───────────────────────────────────────────────────────────

  \path[->]

    % IDLE
    (idle)
      edge [loop left]
        node [left] {\scriptsize $\neg\mathit{req}{\wedge}\neg\mathit{flush}$}
        (idle)
      edge
        node [left] {\scriptsize $\mathit{dc\_req}$}
        (lookup)
      edge
        node [above] {\scriptsize $\mathit{dc\_flush}$}
        (flush)

    % LOOKUP
    (lookup)
      edge [bend right=35]
        node [right] {\scriptsize $\mathit{hit}$}
        (idle)
      edge
        node [left]  {\scriptsize $\neg\mathit{hit}{\wedge}\neg\mathit{dirty}$}
        (miss)
      edge
        node [above] {\scriptsize $\neg\mathit{hit}{\wedge}\mathit{dirty\_victim}$}
        (evict_aw)

    % EVICT_AW
    (evict_aw)
      edge [loop left]
        node [left] {\scriptsize $\neg\mathit{awready}$}
        (evict_aw)
      edge
        node [left] {\scriptsize $\mathit{awready}$}
        (evict_w)

    % EVICT_W
    (evict_w)
      edge [loop left]
        node [left] {\scriptsize $\mathit{wready}{\wedge}\neg\mathit{wlast}$}
        (evict_w)
      edge
        node [left] {\scriptsize $\mathit{wready}{\wedge}\mathit{wlast}$}
        (evict_b)

    % EVICT_B
    (evict_b)
      edge [loop left]
        node [left] {\scriptsize $\neg\mathit{bvalid}$}
        (evict_b)
      edge [bend left=20]
        node [below right] {\scriptsize $\mathit{bvalid}{\wedge}\mathit{ok}{\wedge}\neg\mathit{fm}$}
        (miss)
      edge [out=60, in=220]
        node [below] {\scriptsize $\mathit{bvalid}{\wedge}\mathit{ok}{\wedge}\mathit{fm}{\wedge}\mathit{more}$}
        (flush)
      edge [out=85, in=200]
        node [left] {\scriptsize $\mathit{bvalid}{\wedge}\mathit{ok}{\wedge}\mathit{fm}{\wedge}\mathit{done}$}
        (idle)
      edge
        node [below] {\scriptsize $\mathit{bvalid}{\wedge}\mathit{bresp}{\neq}0$}
        (err)

    % MISS
    (miss)
      edge [loop right]
        node [right] {\scriptsize $\neg\mathit{arready}$}
        (miss)
      edge
        node [left] {\scriptsize $\mathit{arready}$}
        (fill)

    % FILL
    (fill)
      edge [loop left]
        node [left] {\scriptsize $\mathit{rvalid}{\wedge}\neg\mathit{rlast}$}
        (fill)
      edge
        node [left]
        {\scriptsize $\mathit{rvalid}{\wedge}\mathit{rlast}{\wedge}\mathit{ok}$}
        (fill_resp)
      edge [out=0, in=270]
        node [right] {\scriptsize $\mathit{rvalid}{\wedge}\mathit{rresp}{\neq}0$}
        (err)

    % FILL_RESP
    (fill_resp)
      edge [bend right=65]
        node [right] {\scriptsize $\mathbf{1}$}
        (idle)

    % ERR
    (err)
      edge [out=100, in=330]
        node [right] {\scriptsize $\mathbf{1}$}
        (idle)

    % FLUSH
    (flush)
      edge
        node [right] {\scriptsize $\mathit{dirty}$}
        (flush_ep)
      edge [loop right]
        node [right] {\scriptsize $\neg\mathit{dirty}{\wedge}\mathit{more}$}
        (flush)
      edge [bend right=20]
        node [above] {\scriptsize $\neg\mathit{dirty}{\wedge}\mathit{done}$}
        (idle)

    % FLUSH_EVICT_PREP
    (flush_ep)
      edge [bend right=40]
        node [below] {\scriptsize $\mathbf{1}$}
        (evict_aw);

\end{tikzpicture}
}% end resizebox
\caption{Moore-FSM: \texttt{asDCache} Data Cache Controller
  (\texttt{ip\_dcache/src/asDCache.sv}).
  \textbf{Abbreviations on edges:}
  $\mathit{dirty\_victim}$: victim way is valid \emph{and} dirty;
  $\mathit{fm}$ (\texttt{flush\_mode\_r}): eviction was triggered by a
  flush, not a normal miss;
  $\mathit{more}$: flush iterator has not yet scanned all sets/ways;
  $\mathit{done}$: flush complete (all sets/ways iterated);
  $\mathit{ok}$: AXI4 response is OKAY ($\mathit{bresp}{=}0$ or
  $\mathit{rresp}{=}0$);
  $\mathit{wr}$: the CPU request was a store (\texttt{lk\_wr\_r}).
  \textbf{Outputs} $\mathit{rvalid}{=}{\neg}\mathit{wr}$ in
  \textsc{Fill\_Resp}: only a load returns data to the CPU.
  Self-loop label \textbf{1} denotes an unconditional single-cycle
  transition.}
\label{fig:asDCacheFsm01}
\end{figure}

% ─────────────────────────────────────────────────────────────────────────────
\paragraph{State-by-state description of \texttt{asDCache}:}
% ─────────────────────────────────────────────────────────────────────────────

\begin{description}

  \item[\texttt{IDLE}]
    The cache is ready for a new CPU request.
    \texttt{dc\_stall\_o} is de-asserted.
    On \texttt{dc\_flush\_i}: initialise \texttt{flush\_idx\_r} and
    \texttt{flush\_way\_r} to zero, assert stall, go to \textsc{Flush}.
    On \texttt{dc\_req\_i}: latch the full request (address, size, wr,
    wdata, wstrb) into the \texttt{lk\_*\_r} register bank; present
    \texttt{req\_idx\_s} to the data and tag SRAM read ports; assert
    stall; go to \textsc{Lookup}.

  \item[\texttt{LOOKUP}]
    Registered SRAM output (\texttt{data\_rdata\_s}, \texttt{tag\_rdata\_s})
    is valid from the previous-cycle read initiated in \textsc{Idle}.
    Combinatorial valid/dirty arrays (\texttt{vd\_valid\_s},
    \texttt{vd\_dirty\_s}) are read via \texttt{lk\_idx\_r}.
    \begin{itemize}
      \item \textbf{Hit:}
        update PLRU.
        \emph{Read hit} -- extract and sign/zero-extend the double word
        selected by \texttt{lk\_dw\_sel\_r}; assert \texttt{dc\_rdata\_o},
        \texttt{dc\_rvalid\_o}; de-assert stall; return to \textsc{Idle}.
        \emph{Write hit} -- merge write data into the cache line via
        \texttt{merge\_write()}; set dirty bit through the combinatorial
        SRAM write-port; de-assert stall; return to \textsc{Idle}.
        In both cases total latency is \textbf{2 cycles}.
      \item \textbf{Miss, clean victim:}
        no write-back needed; drive AXI4 AR with the line-aligned miss
        address; assert \texttt{ar\_valid\_r}; go to \textsc{Miss}.
      \item \textbf{Miss, dirty victim:}
        write-back required; latch the eviction line
        (\texttt{evict\_line\_r}) from the registered SRAM output; drive
        AXI4 AW channel; assert \texttt{aw\_valid\_r};
        clear \texttt{flush\_mode\_r}; go to \textsc{Evict\_AW}.
    \end{itemize}

  \item[\texttt{EVICT\_AW}]
    AXI4 AW-channel handshake: \texttt{awvalid} held high until
    \texttt{awready} is asserted by the interconnect.
    On acknowledgement: clear \texttt{aw\_valid\_r}; go to
    \textsc{Evict\_W}.

  \item[\texttt{EVICT\_W}]
    AXI4 W-channel data phase: \texttt{wvalid} is driven high
    (combinatorial from state).
    The eviction line (\texttt{evict\_line\_r}, 256\,bit) is streamed in
    four 64-bit beats; \texttt{beat\_r} counts 0\ldots{}3.
    On each \texttt{wready}: increment \texttt{beat\_r} until
    $\mathit{beat\_r}{=}3$ (\texttt{wlast} asserted); go to
    \textsc{Evict\_B}.

  \item[\texttt{EVICT\_B}]
    AXI4 B-channel response: \texttt{bready} held high; FSM waits for
    \texttt{bvalid}.
    On \texttt{bvalid}:
    \begin{itemize}
      \item $\mathit{bresp}{\neq}0$: assert \texttt{dc\_err\_o};
        de-assert stall; go to \textsc{Err}.
      \item $\mathit{bresp}{=}0$ and $\neg\mathit{flush\_mode\_r}$
        (normal miss eviction): proceed to fill the line; drive AR
        channel; go to \textsc{Miss}.
      \item $\mathit{bresp}{=}0$ and $\mathit{flush\_mode\_r}$ and
        more flush work remains (\texttt{flush\_way\_r} or
        \texttt{flush\_idx\_r} not exhausted): invalidate the evicted
        line via the combinatorial write-port; advance the flush
        pointer; go back to \textsc{Flush}.
      \item $\mathit{bresp}{=}0$ and $\mathit{flush\_mode\_r}$ and all
        sets/ways exhausted: assert \texttt{dc\_flush\_done\_o};
        de-assert stall; go to \textsc{Idle}.
    \end{itemize}

  \item[\texttt{MISS}]
    AXI4 AR-channel handshake for a cache-line fill.
    \texttt{arvalid} held high; waits for \texttt{arready}.
    On acknowledgement: clear \texttt{ar\_valid\_r}; go to \textsc{Fill}.

  \item[\texttt{FILL}]
    AXI4 R-channel data phase: \texttt{rready} held high.
    Four 64-bit beats are assembled into \texttt{fill\_buf\_r} /
    \texttt{fill\_line\_s}.
    \begin{itemize}
      \item $\mathit{rresp}{\neq}0$: assert \texttt{dc\_err\_o};
        de-assert stall; go to \textsc{Err}.
      \item $\mathit{rvalid}{\wedge}\neg\mathit{rlast}$: store beat into
        \texttt{fill\_buf\_r}; increment \texttt{beat\_r}; remain in
        \textsc{Fill}.
      \item $\mathit{rvalid}{\wedge}\mathit{rlast}{\wedge}\mathit{ok}$:
        on the last beat the combinatorial write-port decode writes
        data, tag, and valid bit to the SRAMs; if the original request
        was a write-miss (\texttt{lk\_wr\_r}), the write data is merged
        into the fill line before the SRAM write; update PLRU; for a
        read, latch the extracted data word into \texttt{rdata\_r};
        go to \textsc{Fill\_Resp}.
    \end{itemize}

  \item[\texttt{FILL\_RESP}]
    One-cycle response state after a successful fill.
    \texttt{dc\_stall\_o} is de-asserted.
    For a load: assert \texttt{dc\_rvalid\_o} with the latched data word.
    For a write-miss: no rdata output; the merged line has already been
    written to the SRAM in \textsc{Fill}.
    Return to \textsc{Idle}.

  \item[\texttt{ERR}]
    One-cycle error state.
    \texttt{dc\_err\_o} pulsed; \texttt{dc\_stall\_o} de-asserted.
    The affected cache line is \emph{not} written; valid bit is not set.
    Return to \textsc{Idle}.

  \item[\texttt{FLUSH}]
    The flush iterator examines one (set, way) pair per cycle.
    The combinatorial valid/dirty read address is driven by
    \texttt{flush\_idx\_r} (via \texttt{vd\_rd\_addr\_s} mux).
    \begin{itemize}
      \item \textbf{Dirty line:} present \texttt{flush\_idx\_r} to the
        data and tag SRAM read port (one cycle of registered latency
        is needed); go to \textsc{Flush\_Evict\_Prep}.
      \item \textbf{Clean or invalid line:} invalidate the line via the
        combinatorial write-port; advance \texttt{flush\_way\_r} (then
        \texttt{flush\_idx\_r}); remain in \textsc{Flush} until all
        32\,$\times$\,4 = 128 entries are scanned, then assert
        \texttt{dc\_flush\_done\_o} and return to \textsc{Idle}.
    \end{itemize}

  \item[\texttt{FLUSH\_EVICT\_PREP}]
    One-cycle pipeline bubble that absorbs the SRAM read latency.
    In the previous cycle (\textsc{Flush}) the SRAM read address was
    already \texttt{flush\_idx\_r}; the registered output is now valid.
    Latch \texttt{evict\_line\_r} from \texttt{data\_rdata\_s[flush\_way\_r]};
    compute the eviction address from
    \texttt{tag\_rdata\_s[flush\_way\_r]};
    assert \texttt{aw\_valid\_r}; set \texttt{flush\_mode\_r} to
    distinguish this write-back from a normal dirty-victim eviction;
    go to \textsc{Evict\_AW}.

\end{description}

% ─────────────────────────────────────────────────────────────────────────────
\paragraph{Scratchpad Controller (\texttt{asScratch}) -- 2-state FSM.}
% ─────────────────────────────────────────────────────────────────────────────

The \texttt{asScratch} module connects directly to the D-Cache CPU
interface (\texttt{as\_dcache\_if}) and provides a single-port 64-bit
wide SRAM (\texttt{as\_sram\_sp64}) as a byte-addressable scratchpad.
Its controller is a minimal two-state machine whose state is captured by
the single registered bit \texttt{busy\_r} (Fig.~\ref{fig:asScratchFsm01}).

\begin{figure}[H]
\centering
\begin{tikzpicture}[shorten >=1pt, node distance=4.0cm, on grid, auto,
  every state/.style={fill=blue!10}, initial text=]

  % ── State nodes ──────────────────────────────────────────────────────────

  \node[state with output, initial above, accepting] (idle)
    { \texttt{IDLE}
      \nodepart{lower}
      \begin{tiny}$\begin{array}{l}
        \mathit{busy\_r}{=}0 \\
        \mathit{rvalid}{=}0
      \end{array}$\end{tiny}
    };

  \node[state with output] (busy) [right=of idle]
    { \texttt{BUSY}
      \nodepart{lower}
      \begin{tiny}$\begin{array}{l}
        \mathit{busy\_r}{=}1 \\
        \mathit{rvalid}{=}1
      \end{array}$\end{tiny}
    };

  % ── Transitions ───────────────────────────────────────────────────────────

  \path[->]
    (idle)
      edge [loop left]
        node [left] {\scriptsize $\neg\mathit{req}$ or $(\mathit{req}{\wedge}\mathit{wr})$}
        (idle)
      edge [bend left=20]
        node [above] {\scriptsize $\mathit{req}{\wedge}\neg\mathit{wr}$}
        (busy)
    (busy)
      edge [bend left=20]
        node [below] {\scriptsize $\mathbf{1}$}
        (idle);

\end{tikzpicture}
\caption{Moore-FSM: \texttt{asScratch} Scratchpad Controller
  (\texttt{ip\_scratch/src/asScratch.sv}).
  \textbf{Note:} The \texttt{dc\_stall\_o} output is a Mealy signal
  (\textit{stall} $= \mathit{req} \wedge \neg\mathit{wr} \wedge
  \neg\mathit{busy\_r}$) and therefore depends on both state and
  inputs; it is not shown in the state node.
  Stores (\textit{req}~$\wedge$~\textit{wr}) are single-cycle and
  never cause a stall; the SRAM write completes in the same cycle.
  Only loads (\textit{req}~$\wedge$~$\neg$\textit{wr}) require one
  extra cycle for the registered SRAM read output.}
\label{fig:asScratchFsm01}
\end{figure}

\begin{description}

  \item[\texttt{IDLE} ($\mathit{busy\_r}{=}0$)]
    The scratchpad is ready.
    \begin{itemize}
      \item \textbf{No request} ($\neg\mathit{req}$): remain idle; SRAM
        chip-enable (\texttt{cen\_i}) is de-asserted.
      \item \textbf{Store} ($\mathit{req}{\wedge}\mathit{wr}$):
        assert \texttt{cen\_i}, \texttt{we\_i}, and the byte-write
        enables \texttt{wbe\_i} from \texttt{dc\_wstrb\_i}; the
        SRAM completes the write in the current cycle; remain in
        \textsc{Idle} with no stall.
        The \texttt{dc\_stall\_o} output is low for stores.
      \item \textbf{Load} ($\mathit{req}{\wedge}\neg\mathit{wr}$):
        latch \texttt{dc\_size\_i} and \texttt{dc\_addr[2:0]} into
        \texttt{size\_r} and \texttt{boff\_r}; assert \texttt{cen\_i}
        with \texttt{we\_i\,=\,0}; assert \texttt{dc\_stall\_o}
        (Mealy output); set \texttt{busy\_r\,=\,1}; go to \textsc{Busy}.
    \end{itemize}

  \item[\texttt{BUSY} ($\mathit{busy\_r}{=}1$)]
    The SRAM registered read output \texttt{sram\_rdata\_s} is now
    valid.
    The combinatorial output logic sign/zero-extends the 64-bit SRAM
    word according to \texttt{size\_r} (byte, halfword, word, double)
    and byte offset \texttt{boff\_r}, following the RISC-V load
    semantics (\texttt{LB}/\texttt{LH}/\texttt{LW}/\texttt{LD}/
    \texttt{LBU}/\texttt{LHU}/\texttt{LWU}).
    \texttt{dc\_rvalid\_o} is asserted; \texttt{dc\_stall\_o} is
    de-asserted (Mealy: \textit{req} would be 0 because the CPU sees
    stall and has not issued a new request).
    Unconditional transition back to \textsc{Idle}.

\end{description}
